\documentclass[10pt,a4paper]{amsart}
\usepackage{amsmath,amssymb,amsthm,graphicx}
% For commutative diagrams we use Xy-pic
\usepackage[all]{xy}

% We use 2cell for 2-commutative diagrams.
\xyoption{2cell}
% Theorem environments.
%
\theoremstyle{plain}
\newtheorem{theorem}[subsection]{Theorem}
\newtheorem{proposition}[subsection]{Proposition}
\newtheorem{lemma}[subsection]{Lemma}

\theoremstyle{definition}
\newtheorem{definition}[subsection]{Definition}
\newtheorem{example}[subsection]{Example}
\newtheorem{exercise}[subsection]{Exercise}
\newtheorem{situation}[subsection]{Situation}

\theoremstyle{remark}
\newtheorem{remark}[subsection]{Remark}
\newtheorem{remarks}[subsection]{Remarks}

\numberwithin{equation}{subsection}

% Macros
%
\def\lim{\mathop{\mathrm{lim}}\nolimits}
\def\colim{\mathop{\mathrm{colim}}\nolimits}
\def\Spec{\mathop{\mathrm{Spec}}}
\def\Hom{\mathop{\mathrm{Hom}}\nolimits}
\def\Ext{\mathop{\mathrm{Ext}}\nolimits}
\def\SheafHom{\mathop{\mathcal{H}\!\mathit{om}}\nolimits}
\def\SheafExt{\mathop{\mathcal{E}\!\mathit{xt}}\nolimits}
\def\Sch{\mathit{Sch}}
\def\Mor{\mathop{\mathrm{Mor}}\nolimits}
\def\Ob{\mathop{\mathrm{Ob}}\nolimits}
\def\Sh{\mathop{\mathit{Sh}}\nolimits}
\def\NL{\mathop{N\!L}\nolimits}
\def\CH{\mathop{\mathrm{CH}}\nolimits}
\def\proetale{{pro\text{-}\acute{e}tale}}
\def\etale{{\acute{e}tale}}
\def\QCoh{\mathit{QCoh}}
\def\Ker{\mathop{\mathrm{Ker}}}
\def\Im{\mathop{\mathrm{Im}}}
\def\Coker{\mathop{\mathrm{Coker}}}
\def\Coim{\mathop{\mathrm{Coim}}}

% Boxtimes
%
\DeclareMathSymbol{\boxtimes}{\mathbin}{AMSa}{"02}

%
% Macros for moduli stacks/spaces
%
\def\QCohstack{\mathcal{QC}\!\mathit{oh}}
\def\Cohstack{\mathcal{C}\!\mathit{oh}}
\def\Spacesstack{\mathcal{S}\!\mathit{paces}}
\def\Quotfunctor{\mathrm{Quot}}
\def\Hilbfunctor{\mathrm{Hilb}}
\def\Curvesstack{\mathcal{C}\!\mathit{urves}}
\def\Polarizedstack{\mathcal{P}\!\mathit{olarized}}
\def\Complexesstack{\mathcal{C}\!\mathit{omplexes}}
% \Pic is the operator that assigns to X its picard group, usage \Pic(X)
% \Picardstack_{X/B} denotes the Picard stack of X over B
% \Picardfunctor_{X/B} denotes the Picard functor of X over B
\def\Pic{\mathop{\mathrm{Pic}}\nolimits}
\def\Picardstack{\mathcal{P}\!\mathit{ic}}
\def\Picardfunctor{\mathrm{Pic}}
\def\Deformationcategory{\mathcal{D}\!\mathit{ef}}
\usepackage{cleveref}
\crefname{originaltheorem}{original theorem}{original theorems}
\Crefname{originaltheorem}{Original theorem}{Original theorems}
\makeatletter
\newlabel{theorem-colimit}{{09YQ}{source}{}{sourcepointer-1}{}}
\newlabel{theorem-colimit@cref}{{[originaltheorem][0][]09YQ}{[1][0][]source}{}{}{}}
\makeatother
\crefname{originallemma}{original lemma}{original lemmas}
\Crefname{originallemma}{Original lemma}{Original lemmas}
\makeatletter
\newlabel{algebra-lemma-etale}{{00U2}{source}{}{sourcepointer-2}{}}
\newlabel{algebra-lemma-etale@cref}{{[originallemma][0][]00U2}{[1][0][]source}{}{}{}}
\makeatother
\crefname{originaltheorem}{original theorem}{original theorems}
\Crefname{originaltheorem}{Original theorem}{Original theorems}
\makeatletter
\newlabel{theorem-cech-ss}{{03OW}{source}{}{sourcepointer-3}{}}
\newlabel{theorem-cech-ss@cref}{{[originaltheorem][0][]03OW}{[1][0][]source}{}{}{}}
\makeatother
\title{Cohomology of the base change along a directed colimit of algebras as the colimit of the cohomologies}
\author{The Stacks Project Authors}
\date{stacks-project:etale-cohomology.tex:local-snapshot-20261011; GNU FDL 1.2 or later}
\begin{document}
\maketitle
\noindent\textit{Editorial note.} All delivered mathematical text is the author's own, copied from the pinned Stacks Project chapter file etale-cohomology.tex; no statement or proof step was rewritten, reordered or removed.\par\medskip
\noindent\textit{Editorial note.} The proof is the author's complete argument; it identifies the limit presentation and applies the cited theorem.\par\medskip
\noindent\textit{Editorial note.} Every labelled reference occurring inside the counted spans is carried in assembly.external\_references as a non-load-bearing navigation pointer (same-chapter ones at their definition site, cross-chapter ones at the invocation inside the counted span), so no load-bearing mathematics is replaced by a stub. No bibliography citation occurs inside the counted spans.\par\medskip
\noindent\textit{Editorial note.} The checked chapter list for same-name collisions is the set of stacks chapters staged in this repository: chow, cohomology, divisors, etale-cohomology, flat, more-algebra, more-morphisms, morphisms, spaces-more-morphisms, varieties together with this card's chapter etale-cohomology.\par\medskip
\noindent\textit{Editorial note.} The preamble is the chapter's own native preamble excerpt already used for this source: the Xy-pic setup (native preamble lines 14-18) and the native theorem declarations with the mathematical macros (native preamble lines 154-217).\par\medskip
\noindent\textit{Allowed standard prerequisites (editorial).}\begin{itemize}
\item Base change of a scheme along a morphism of affine schemes and the resulting cover. Same chapter, the description of Y\_i = X x\_\{Spec(A)\} Spec(B\_i) and of the pullbacks G\_i of the sheaf F. Supplies the schemes and sheaves appearing in the asserted identification.
\item Commutation of etale cohomology with directed colimits of affine base schemes. Same chapter, the theorem stating that etale cohomology of a quasi-compact quasi-separated scheme over a limit of affine schemes is the colimit of the cohomologies, of which this lemma is a special case. Supplies the result applied in the proof.
\end{itemize}\medskip
\ifdefined\hypertarget\hypertarget{sourcepointer-1}{}\fi
\noindent\textit{Source reference.} Original theorem 09YQ: sources/stacks-etale-cohomology/pinned-original-source.tex, source lines 7400--7400. Navigation pointer invoked at proof line 7438; the counted span contains the pointer only and the cited result is not reproduced inside it, so the pointer is registered as a non-load-bearing navigation stub in the established form for Stacks material. No bibliography citation is involved.\par\medskip
\ifdefined\hypertarget\hypertarget{sourcepointer-2}{}\fi
\noindent\textit{Source reference.} Original lemma 00U2: sources/stacks-etale-cohomology/pinned-original-source.tex, source lines 7490--7490. Navigation pointer invoked at proof line 7490; the counted span contains the pointer only and the cited result is not reproduced inside it, so the pointer is registered as a non-load-bearing navigation stub in the established form for Stacks material. No bibliography citation is involved.\par\medskip
\ifdefined\hypertarget\hypertarget{sourcepointer-3}{}\fi
\noindent\textit{Source reference.} Original theorem 03OW: sources/stacks-etale-cohomology/pinned-original-source.tex, source lines 1771--1771. Navigation pointer invoked at proof line 7456; the counted span contains the pointer only and the cited result is not reproduced inside it, so the pointer is registered as a non-load-bearing navigation stub in the established form for Stacks material. No bibliography citation is involved.\par\medskip
\begin{lemma}
\label{lemma-directed-colimit-cohomology}
Let $A$ be a ring, $(I, \leq)$ a directed set and $(B_i, \varphi_{ij})$ a
system of $A$-algebras. Set $B = \colim_{i\in I} B_i$. Let $X \to \Spec(A)$
be a quasi-compact and quasi-separated morphism of schemes. Let
$\mathcal{F}$ an abelian sheaf on $X_\etale$.
Denote $Y_i = X \times_{\Spec(A)} \Spec(B_i)$,
$Y = X \times_{\Spec(A)} \Spec(B)$,
$\mathcal{G}_i = (Y_i \to X)^{-1}\mathcal{F}$ and
$\mathcal{G} = (Y \to X)^{-1}\mathcal{F}$. Then
$$
H_\etale^p(Y, \mathcal{G}) =
\colim_{i\in I} H_\etale^p (Y_i, \mathcal{G}_i).
$$
\end{lemma}

\begin{proof}
This is a special case of Theorem \ref{theorem-colimit}.
We also outline a direct proof as follows.
\begin{enumerate}
\item Given $V \to Y$ \'etale with $V$ quasi-compact and
quasi-separated, there exist $i\in I$ and $V_i \to Y_i$ such that
$V = V_i \times_{Y_i} Y$.
If all the schemes considered were affine, this would correspond to the
following algebra statement: if $B = \colim B_i$ and $B \to C$ is \'etale,
then there exist $i \in I$ and $B_i \to C_i$ \'etale such that
$C \cong B \otimes_{B_i} C_i$.
This is proved in Algebra, Lemma \ref{algebra-lemma-etale}.
\item In the situation of (1) show that
$\mathcal{G}(V) = \colim_{i' \geq i} \mathcal{G}_{i'}(V_{i'})$
where $V_{i'}$ is the base change of $V_i$ to $Y_{i'}$.
\item By (1), we see that for every \'etale covering
$\mathcal{V} = \{V_j \to Y\}_{j\in J}$ with $J$ finite and the
$V_j$s quasi-compact and quasi-separated, there exists $i \in I$ and
an \'etale covering $\mathcal{V}_i = \{V_{ij} \to Y_i\}_{j \in J}$
such that $\mathcal{V} \cong \mathcal{V}_i \times_{Y_i} Y$.
\item Show that (2) and (3) imply
$$
\check H^*(\mathcal{V}, \mathcal{G})=
\colim_{i\in I} \check H^*(\mathcal{V}_i, \mathcal{G}_i).
$$
\item Cleverly use the {\v C}ech-to-cohomology spectral sequence
(Theorem \ref{theorem-cech-ss}).
\end{enumerate}
\end{proof}

\end{document}
